% =====================================================================================================
\chapter*{Abstract}\addcontentsline{toc}{chapter}{Abstract}
% =====================================================================================================
Can a single X-ray observation tell whether an accreting compact object is a black hole (BH) or a neutron star (NS), for a source
that the classifier has never seen? The XRB-pilot project (versions v1--v14) answered this question with a strict design:
dynamically confirmed BHs and burst- or pulse-confirmed NSs only; leave-one-source-out evaluation; source-level metrics with a
class-stratified source bootstrap; one or a few pre-registered primary comparisons per version, judged by whether the 95\%
interval excludes zero; and frozen models, rules and pipelines for every test on new data.

Using RXTE/PCA spectra of 31 sources (7 BH, 24 NS), a single 5--25\,keV spectrum classifies unseen sources well above chance
(source balanced accuracy 0.80--0.87, source AUC 0.88--0.95), but essentially all of the information is in two count-space
hardness ratios: the full spectral shape, intensity, 3--5 and 25--60\,keV bands, flexible algorithms with nested selection, 17 times
more observations per source and 15 additional sources add nothing detectable to the ranking of unseen sources. Variability-defined states show why: the colour
classifier separates BHs from NSs in soft states (observation-level AUC 0.954) but not in hard states (0.646; difference
$+0.31$ [$+0.10$, $+0.51$]). Inside the hard state, 0.016--4\,Hz variability measured with cross-detector cospectra adds
information: $+0.26$ [$+0.05$, $+0.46$] in AUC on the RXTE sample, consistently positive on NICER, $+0.15$ [$+0.03$, $+0.26$] when
RXTE and NICER are pooled (86 sources, 16 dynamically confirmed BHs; not an independent confirmation), and $+0.18$
[$+0.05$, $+0.34$] in a prospective test with a frozen pipeline and a pre-chosen random forest on new NICER observations. At fixed
colour the characteristic frequency of BH variability is lower by 0.3--0.6\,dex in every data set but one, with pooled $p=0.0012$
but non-significant single prospective tests. A burst catalogue, persistent BHs, event-mode frequencies up to 1\,kHz and MAXI/GSC
light curves (including an exact reproduction of a published classifier) qualify these results. All prospective confirmations are
at the observation level on the same 16 BHs; a source-level test on new BH transients is the necessary next step. This report
documents every phase, the full mathematics of features, models and statistics, and the mapping from the archived data to the
machine-learning matrices; it re-runs no analysis and every number is read from files committed in the repository.
